%% 02-lineage-and-state-of-the-art.tex -- generated by tools/build_from_docx.py; edit freely

\hypertarget{lineage-and-state-of-the-art}{%
\section{Lineage and State of the Art}\label{lineage-and-state-of-the-art}}

The four-capability structure at the center of this framework was not
invented for this project. It arrived from earlier work on the
management of fifth-generation networks, where it was devised to address
a problem that had become unavoidable: the management systems of a
network were failing to scale with the network itself.

Being explicit about that inheritance matters, because several positions
this framework takes look like preferences and are in fact consequences.
A reviewer who proposes a central arbiter over the decision layer is not
proposing a variation on this framework; they are proposing the
architecture the pattern was created to replace.

\hypertarget{the-problem-that-produced-the-pattern}{%
\subsection{The Problem that Produced the Pattern}\label{the-problem-that-produced-the-pattern}}

Conventional management gathers information centrally, reasons about it
centrally, and issues instructions outward. Each of those steps becomes
more expensive as the managed system grows, and the growth in question
was not incremental. Network slicing multiplied the number of
independently managed entities, virtualization multiplied the things
that could change, and the timescales on which decisions were expected
shortened at the same time. A single management system observing
everything and deciding everything was on a trajectory that ended in it
becoming the bottleneck of the network it was meant to be optimizing.

The response was to stop treating management as something applied to the
network from outside, and to embed it within the managed entities
themselves. Each entity would carry its own observation, its own
analysis, its own decision making and its own means of acting: a
complete closed loop, local to the thing it managed, operating at its
own timescale.

That decomposition is the origin of the four capabilities this framework
describes. The prior work, the MonB5G project under Horizon 2020 (grant
agreement 871780), defined them as a Monitoring System, an Analytics
Engine, a Decision Engine and an Actuator, composing a closed management
loop distributable across managed network components, and it established
that distributed artificial-intelligence-driven closed loops are viable
and scalable at the scale of network slicing.

\hypertarget{what-is-inherited-rather-than-chosen}{%
\subsection{What is Inherited Rather than Chosen}\label{what-is-inherited-rather-than-chosen}}

Four positions in this framework follow directly from that origin.

There is no global decision engine. The pattern exists because
centralization did not scale, and reintroducing a central brain to
coordinate the decentralized loops would undo the reason for having
them.

Raw high-rate observation stays in the domain that produced it, while
summaries, features and derived products travel. The traffic saved by
not moving raw telemetry centrally was part of what the decomposition
was for, and the scalability argument is defeated if all observation
converges anyway.

The relationship between management levels is delegation rather than
command. Higher layers set objectives and bounds; they do not steer
equipment on the fast loop. A hierarchy in which the upper layer issued
binding instructions would be a central controller with extra steps.

And the attention this framework pays to where its guardrails cost time,
formalized as the cost classes of Chapter 12, is inherited from the same
source. A guardrail placed on the per-decision path partially
re-centralizes control, so cost class is normative rather than advisory.
A pattern adopted to relieve a scalability bottleneck cannot be
implemented in a way that creates a new one. Table 2-1 states each
inherited position against the argument it follows from.

\begin{wptable}{2-1}

\begin{longtable}[]{@{}
  >{\raggedright\arraybackslash}p{(\columnwidth - 4\tabcolsep) * \real{0.2750}}
  >{\raggedright\arraybackslash}p{(\columnwidth - 4\tabcolsep) * \real{0.4750}}
  >{\raggedright\arraybackslash}p{(\columnwidth - 4\tabcolsep) * \real{0.2501}}@{}}
\caption{Inherited positions and what each follows from}\tabularnewline
\toprule\noalign{}
\begin{minipage}[b]{\linewidth}\raggedright
\textbf{Inherited position}
\end{minipage} & \begin{minipage}[b]{\linewidth}\raggedright
\textbf{What it follows from}
\end{minipage} & \begin{minipage}[b]{\linewidth}\raggedright
\textbf{Where it is specified}
\end{minipage} \\
\midrule\noalign{}
\endfirsthead
\toprule\noalign{}
\begin{minipage}[b]{\linewidth}\raggedright
\textbf{Inherited position}
\end{minipage} & \begin{minipage}[b]{\linewidth}\raggedright
\textbf{What it follows from}
\end{minipage} & \begin{minipage}[b]{\linewidth}\raggedright
\textbf{Where it is specified}
\end{minipage} \\
\midrule\noalign{}
\endhead
\bottomrule\noalign{}
\endlastfoot
\textbf{No global decision engine} & Centralization did not scale; a
central coordinator over the decentralized loops would restore the
bottleneck the pattern was adopted to remove. & Chapters 6 and 10 \\
\textbf{Raw high-rate observation stays in its domain} & The traffic
saved by not moving raw telemetry centrally is part of the scalability
case; summaries, features and derived products travel instead. &
Chapters 4 and 8 \\
\textbf{Delegation, not command, between levels} & An upper layer
issuing binding instructions on the fast loop is a central controller
under another name; it sets objectives and bounds instead. & Sections
3.5.21 and 3.9.2; Chapter 10 \\
\textbf{Guardrail cost is normative} & A guardrail on the per-decision
path partially re-centralizes control, so a pattern adopted to relieve a
bottleneck cannot be implemented in a way that creates one. & Chapter
12 \\
\end{longtable}

\end{wptable}

\hypertarget{the-extension-from-tuple-to-agent}{%
\subsection{The Extension: From Tuple to Agent}\label{the-extension-from-tuple-to-agent}}

What this project adds is the step from a management tuple to an agent,
and it is a larger step than the vocabulary suggests.

In its original form the tuple was a structure for organizing management
functions. The decision capability held logic that had been designed,
reviewed and deployed by the party responsible for the entity it
managed, and the safety of the arrangement rested on that logic being
sound. Boundaries could be assumed, because a trained model has an
implicitly bounded output space and the party who trained it was the
party deploying it.

That arrangement does not survive the conditions this project is
building for. The intelligence occupying these capabilities arrives from
partners whose internals the platform does not inspect, addressing
objectives declared at instantiation rather than at design time, using
techniques that range from trained models through optimizers and digital
twins to language-model agents reasoning in natural language.

An agent in this framework is therefore three things the tuple was not.
It is instantiated against a declared objective rather than merely wired
as a pipeline. It carries declared authority, control scope and failure
behavior. And its reasoning stages are open as to technique.

\hypertarget{technique-neutrality}{%
\subsubsection{Technique neutrality}\label{technique-neutrality}}

An Analytics Engine or a Decision Engine may be realized by classical
machine learning, whether supervised, reinforcement or statistical,
which is the inherited case. It may be realized by language-model or
foundation-model reasoning, including tool use and multi-step planning.
It may be deterministic or rule-based logic, an optimization solver, or
digital-twin evaluation. Or it may be any hybrid of these.

The framework holds identically across all of them, and this is what
AI-native means in this document: not that every component contains a
model, but that the architecture assumes autonomous AI-driven reasoning
as the normal case and is built to contain it safely.

\hypertarget{why-the-role-contracts-are-load-bearing}{%
\subsubsection{Why the role contracts are load-bearing}\label{why-the-role-contracts-are-load-bearing}}

Technique neutrality is what forces everything else. A language-model
agent has no bounded output space by construction; no inspection of it
will produce one, and no assurance from its contributor can substitute
for one.

If such a technique is to hold authority over a live network, and the
project's ambition requires that it can, then the bounding must come
from outside the technique. Constraining Decision Engine output to the
registered action space is what makes a generative Decision Engine
admissible at all. The prohibition on self-approval contains reasoning
whose internal justification cannot be audited line by line. Structured
rationale is a provenance floor precisely because interpretability
cannot be demanded uniformly across such heterogeneous techniques. And
the four-outcome actuation vocabulary with its mandatory read-back means
a confidently worded but wrong decision still cannot masquerade as a
successful change.

These are not hygiene applied to a system that would mostly work without
them. They are the conditions under which an unbounded technique can be
given real authority, which makes the guardrails the enabling condition
for agentic operation rather than an overlay on it.

\hypertarget{consequence-for-the-model-domain}{%
\subsubsection{Consequence for the model domain}\label{consequence-for-the-model-domain}}

Because technique is open, the platform cannot treat models as someone
else's concern. Model lifecycle, meaning training, retraining,
versioning, serving, evaluation, promotion and withdrawal, together with
the data pipelines that feed it, is part of the architecture, with
defined relationships to the roles that consume it.

\hypertarget{what-this-task-claims-as-new}{%
\subsection{What This Task Claims as New}\label{what-this-task-claims-as-new}}

The prior work established that distributed closed loops are viable and
scalable. It did not have to answer what arises when many independently
authored loops, contributed by different partners, share the same
transport and the same infrastructure. That question is this framework's
subject, and four contributions follow from it.

The first is multi-tenant admission. The original pattern assumed the
party deploying a loop was the party that built it. This framework is a
platform on which independent contributors are admitted, at authorities
the platform determines, with their declarations checked rather than
accepted. It is flexible in what may be declared and strict in requiring
that it be declared and in holding the contribution to it.

The second is mediation between independently reasoning loops. The
inherited pattern distributed decision making but did not address what
happens when distributed decisions interact. Where two agents reason
correctly in their own domains and their actions conflict, the
information needed to see the conflict exists in neither of them. The
architecture work names this gap directly, observing that agentic
conflict handling lacks operational guardrails for safe, deterministic
enforcement in real networks. This framework's contribution is the
agent-side contract and the guardrails that make external mediation
load-bearing rather than advisory.

The third is verified actuation and honest assurance. That an
acknowledgement is not proof of change, that an unknown outcome is a
state to be resolved rather than a reason to retry, and that unverified
and unresolved actions appear in assurance in their own right rather
than being absorbed into a success figure. This discipline is what
allows the platform to report on outcomes rather than on activity.

The fourth is the extension of the whole construct to sixth-generation
wireless transport, where the managed resource is the radio link itself
rather than only a virtualized slice, together with the recursive
application of the tuple to the management hierarchy, so that
coordination between domains is governed by the same disciplines as
action within them. Table 2-2 states the four contributions against what
the inherited pattern supplied.

\begin{wptable}{2-2}

\begin{longtable}[]{@{}
  >{\raggedright\arraybackslash}p{(\columnwidth - 4\tabcolsep) * \real{0.2501}}
  >{\raggedright\arraybackslash}p{(\columnwidth - 4\tabcolsep) * \real{0.3375}}
  >{\raggedright\arraybackslash}p{(\columnwidth - 4\tabcolsep) * \real{0.4125}}@{}}
\caption{Contributions of this task beyond the inherited
pattern}\tabularnewline
\toprule\noalign{}
\begin{minipage}[b]{\linewidth}\raggedright
\textbf{Contribution}
\end{minipage} & \begin{minipage}[b]{\linewidth}\raggedright
\textbf{What the inherited pattern supplied}
\end{minipage} & \begin{minipage}[b]{\linewidth}\raggedright
\textbf{What this framework adds}
\end{minipage} \\
\midrule\noalign{}
\endfirsthead
\toprule\noalign{}
\begin{minipage}[b]{\linewidth}\raggedright
\textbf{Contribution}
\end{minipage} & \begin{minipage}[b]{\linewidth}\raggedright
\textbf{What the inherited pattern supplied}
\end{minipage} & \begin{minipage}[b]{\linewidth}\raggedright
\textbf{What this framework adds}
\end{minipage} \\
\midrule\noalign{}
\endhead
\bottomrule\noalign{}
\endlastfoot
\textbf{Multi-tenant admission} & Each loop built and deployed by one
party, its logic trusted because its author was its operator. &
Independent contributors admitted at an authority the platform
determines, with their declarations checked rather than accepted
(Chapters 3 and 9). \\
\textbf{Mediation between independently reasoning loops} & Decision
making distributed across entities, with no account of distributed
decisions that interact. & The agent-side contract with the conflict
fabric and the guardrails that make external mediation load-bearing: no
self-approval, a fixed verdict vocabulary and a token checked at
actuation (Chapters 6 and 10). \\
\textbf{Verified actuation and honest assurance} & Actuation as the last
step of the loop, its success taken from the
controller\textquotesingle s acknowledgement. & Four typed outcomes,
independent read-back, an unknown outcome treated as a state to resolve,
and assurance that reports unverified and unresolved actions in their
own right (Chapters 7 and 9). \\
\textbf{Sixth-generation transport and recursive hierarchy} & Loops over
virtualized slices within a single management level. & The radio link as
a managed resource, and the tuple applied recursively to the IDMO--DMO
hierarchy, so that coordination between domains obeys the same
disciplines as action within them (Chapters 3 and 10). \\
\end{longtable}

\end{wptable}

\hypertarget{relation-to-existing-practice}{%
\subsection{Relation to Existing Practice}\label{relation-to-existing-practice}}

This framework is deliberately confined to the practice of the relevant
standards bodies, and it is worth stating where it sits relative to
each.

The closed-loop automation levels and the intent-driven management
interfaces of the third-generation partnership project supply the
management vocabulary this framework binds to rather than restates. The
operational states of the service lifecycle are adopted directly, as is
the service-level assurance framing.

The open radio access network architecture supplies the timescale
separation this framework adopts for its latency classes, and the
policy-guidance pattern between the non-real-time and near-real-time
controllers is the precedent for the action envelopes of Section 3.9.2.
The framework's actuation bindings target the standard interfaces of
that architecture rather than inventing southbound protocols.

Where this framework goes beyond both is in the multi-tenancy of the
intelligence itself. Neither body currently specifies how an operator
admits a decision-making component built by a third party, bounds what
it may reach, mediates between several such components acting at once,
and reports honestly on whether their actions took effect. Those are the
questions this framework answers, and they are answered in the
vocabulary of those standards rather than alongside it.

\hypertarget{the-standing-of-this-framework}{%
\subsection{The Standing of This Framework}\label{the-standing-of-this-framework}}

The conceptual architecture this project is building does not require
that every managed service be realized this way. Services defined
conventionally, without tuples and without the agentic construct, are
legitimate within it and coexist with agent-based services in the same
deployment. This framework makes no claim that the conventional approach
is inadequate for what it already does well.

The claim is about capability. The agentic construct is what the
architecture requires in order to be AI-native in substance rather than
in label, and to host the concepts that distinguish a sixth-generation
system: autonomous loops operating at machine timescales, contributed
intelligence of arbitrary technique holding real authority over live
infrastructure, coordination between agents that reason independently,
and assurance that speaks to outcomes rather than to actions attempted.

These are not reachable by extending conventional service management.
The architecture may adopt the construct where those capabilities are
wanted, and adopting it entails a good deal more than naming four
capabilities and hoping the intelligence placed in them behaves.
